\documentclass[11pt]{article}
\usepackage[margin=1in]{geometry}
\usepackage{amsmath}
\usepackage{amssymb}
\usepackage{tcolorbox}

\title{Why Does the Inequality Bind? \\ (Understanding $\lambda_1 = p_1 - p_2$)}
\author{ECO 310 - The Deeper Question}
\date{February 23, 2026}

\begin{document}

\maketitle

\section{The Question}

We have: $-p_1 + \lambda_1 + \lambda_3 \leq 0$

After substituting $\lambda_3 = p_2$, this becomes: $\lambda_1 \leq p_1 - p_2$

Why do we assume this holds with \textbf{equality}: $\lambda_1 = p_1 - p_2$?

\section{Short Answer}

\begin{tcolorbox}[colback=yellow!10!white,colframe=orange!75!black,title=Quick Answer]
Actually, we don't always assume it binds!

The condition $\lambda_1 \leq p_1 - p_2$ tells us the \textbf{maximum possible value} of $\lambda_1$.

For perfect substitutes specifically, it turns out that $\lambda_1$ takes its maximum value at the optimum. But this is a result we derive, not something we assume.
\end{tcolorbox}

\section{The Full Story}

\subsection{What KKT Conditions Tell Us}

The Karush-Kuhn-Tucker (KKT) optimality conditions for our problem include:

\begin{align}
\frac{\partial \mathcal{L}}{\partial x_1} &\leq 0, \quad x_1 \geq 0, \quad x_1 \cdot \frac{\partial \mathcal{L}}{\partial x_1} = 0 \\
\frac{\partial \mathcal{L}}{\partial x_2} &\leq 0, \quad x_2 \geq 0, \quad x_2 \cdot \frac{\partial \mathcal{L}}{\partial x_2} = 0
\end{align}

These conditions are \textbf{necessary} for optimality (must be true at any optimal solution).

\subsection{What We Actually Know}

From the guess that $x_1 = 0$ and $x_2 > 0$:

\textbf{We know for certain:}
\begin{itemize}
\item $\lambda_2 = 0$ (from complementary slackness with $x_2 > 0$)
\item $\lambda_3 = p_2$ (from the $x_2$ FOC being exactly zero)
\item $x_2 = u^*$ (from utility constraint with $x_1 = 0$)
\end{itemize}

\textbf{We know only a constraint on $\lambda_1$:}
\begin{itemize}
\item $\lambda_1 \leq p_1 - p_2$ (from the $x_1$ FOC inequality)
\item $\lambda_1 \geq 0$ (all Lagrange multipliers must be non-negative)
\end{itemize}

So $\lambda_1 \in [0, p_1 - p_2]$ is feasible.

\section{Why We Set $\lambda_1 = p_1 - p_2$}

There are a few ways to think about this:

\subsection{Approach 1: It's the Unique Value That Works}

For perfect substitutes, there's a unique optimal solution. If the optimal $\lambda_1$ were strictly less than $p_1 - p_2$, this would suggest the solution is not fully determined by the first-order conditions alone, which contradicts the uniqueness of the corner solution.

\begin{tcolorbox}[colback=blue!5!white,colframe=blue!75!black,title=Intuition]
Think of it this way: the system of optimality conditions should uniquely determine all the variables.

If $\lambda_1$ could be anywhere in $[0, p_1 - p_2]$, we'd have infinitely many "solutions," which doesn't make sense for a unique optimization problem.

The fact that we're at a corner solution (buying only one good) suggests the constraints are "tight" - there's no wiggle room.
\end{tcolorbox}

\subsection{Approach 2: Economic Interpretation}

Remember what $\lambda_1$ means: the shadow price of the constraint $x_1 \geq 0$.

\textbf{Question:} How much would your cost decrease if you could buy $x_1 = -\epsilon$ (slightly negative)?

\textbf{Answer:}
\begin{itemize}
\item You'd "sell" $\epsilon$ units of good 1 and get $p_1 \epsilon$ dollars back
\item You'd need to buy $\epsilon$ more of good 2 to maintain utility, costing $p_2 \epsilon$
\item Net savings: $(p_1 - p_2) \epsilon$ per unit
\end{itemize}

So the shadow price should be exactly $p_1 - p_2$, not something smaller!

\subsection{Approach 3: The Math is Actually More Subtle}

Here's the truth: we don't \textit{assume} it binds. Instead:

\begin{enumerate}
\item We \textbf{guess} that $x_1 = 0$, $x_2 = u^*$ is optimal
\item We derive that IF this is optimal, THEN $\lambda_1 = p_1 - p_2$
\item We check that this solution satisfies ALL the KKT conditions
\item If it does, we've found an optimal solution
\end{enumerate}

\section{The Detailed Argument}

Let me show you why $\lambda_1 = p_1 - p_2$ is the "right" value.

\subsection{Setting Up the Problem}

We want to minimize $-\mathcal{L} = p_1 x_1 + p_2 x_2$ (the actual cost) subject to all constraints.

The Lagrangian with Convention 2 is:
$$\mathcal{L} = -p_1 x_1 - p_2 x_2 + \lambda_1 x_1 + \lambda_2 x_2 + \lambda_3(x_1 + x_2 - u^*)$$

At a maximum of $\mathcal{L}$ (which corresponds to minimum cost):
\begin{itemize}
\item We can't increase $\mathcal{L}$ by changing $x_2$ (since $x_2 > 0$, the derivative must be exactly 0)
\item We can't increase $\mathcal{L}$ by increasing $x_1$ from 0 (derivative must be $\leq 0$)
\item We can't increase $\mathcal{L}$ by changing $\lambda_3$ (derivative must be 0)
\end{itemize}

\subsection{Why the Inequality is Actually Tight}

Consider what happens if $\lambda_1 < p_1 - p_2$ (strict inequality).

Then: $-p_1 + \lambda_1 + \lambda_3 < 0$

This means $\frac{\partial \mathcal{L}}{\partial x_1} < 0$ (strictly negative).

\textbf{Interpretation:} Increasing $x_1$ would \textit{decrease} the Lagrangian.

But wait - we're at $x_1 = 0$. We can't decrease $x_1$ further (nonnegativity constraint).

However, the strict inequality $\frac{\partial \mathcal{L}}{\partial x_1} < 0$ means:
\begin{quote}
"If we could buy less than zero of good 1, the Lagrangian would increase (equivalently, cost would decrease)."
\end{quote}

The shadow price $\lambda_1$ measures exactly how much the Lagrangian would increase per unit if we could violate $x_1 \geq 0$.

From our economic argument: violating $x_1 \geq 0$ by one unit (setting $x_1 = -1$) and adjusting $x_2$ to maintain utility would decrease cost by exactly $p_1 - p_2$.

Therefore: $\lambda_1 = p_1 - p_2$ (the full amount, not something smaller).

\section{What Happens in Other Problems?}

In some optimization problems, you \textit{can} have slack in the inequality (i.e., $\lambda_1 < p_1 - p_2$).

\textbf{Example where it doesn't bind:}

Suppose we had a different utility function (not perfect substitutes), and at the optimum:
\begin{itemize}
\item $x_1 = 0$ (still at the boundary)
\item But switching to $x_1 = -\epsilon$ wouldn't actually help much (maybe the utility loss is too high)
\end{itemize}

In that case, $\lambda_1$ could be less than $p_1 - p_2$.

\begin{tcolorbox}[colback=green!5!white,colframe=green!75!black,title=The Key Difference]
For \textbf{perfect substitutes}: goods are exactly equivalent for utility, so switching between them is purely about prices. This makes the shadow price exactly equal to the price difference.

For \textbf{other utility functions}: goods are not perfect substitutes, so the shadow price depends on marginal utilities, not just prices.
\end{tcolorbox}

\section{Practical Advice for Problem 2b}

\begin{tcolorbox}[colback=yellow!10!white,colframe=orange!75!black,title=What to Write on Your Homework]

\textbf{Approach 1 (Simpler):} Just set $\lambda_1 = p_1 - p_2$ and verify it satisfies all conditions.

\textbf{Approach 2 (More rigorous):}
\begin{enumerate}
\item Derive that $\lambda_1 \leq p_1 - p_2$
\item Use economic reasoning: the shadow price of $x_1 \geq 0$ equals the price difference
\item Conclude $\lambda_1 = p_1 - p_2$
\item Verify $\lambda_1 \geq 0$ requires $p_1 \geq p_2$
\end{enumerate}

Either approach is fine for this problem!
\end{tcolorbox}

\section{Summary}

\begin{tcolorbox}[colback=blue!5!white,colframe=blue!75!black,title=The Bottom Line]

\textbf{Why does $\lambda_1 = p_1 - p_2$ instead of $\lambda_1 < p_1 - p_2$?}

\begin{enumerate}
\item For perfect substitutes, the corner solution is unique and fully determined
\item The shadow price $\lambda_1$ measures the cost savings from relaxing $x_1 \geq 0$
\item This cost savings is exactly $p_1 - p_2$ (the price difference)
\item Therefore $\lambda_1$ takes its maximum feasible value

\item Mathematically: the inequality allows $\lambda_1 \in [0, p_1 - p_2]$, but the economics pins down $\lambda_1 = p_1 - p_2$
\end{enumerate}

\textbf{In general:} Not all inequalities in KKT conditions bind. But for this specific problem with perfect substitutes, it does.
\end{tcolorbox}

\section{Alternative Perspective: Solving Directly}

Actually, there's another way to see this that avoids the question entirely.

We have 5 unknowns: $x_1, x_2, \lambda_1, \lambda_2, \lambda_3$

From our guess and complementary slackness:
\begin{align}
x_1 &= 0 \quad \text{(given)} \\
x_2 &= u^* \quad \text{(from utility constraint)} \\
\lambda_2 &= 0 \quad \text{(from complementary slackness)} \\
\lambda_3 &= p_2 \quad \text{(from FOC for $x_2$)}
\end{align}

That's 4 equations. We have one more FOC inequality:
$$-p_1 + \lambda_1 + \lambda_3 \leq 0$$

Substituting $\lambda_3 = p_2$:
$$\lambda_1 \leq p_1 - p_2$$

Now, we want to find THE optimal solution. Since the problem has a unique minimum (perfect substitutes have unique corner solutions), all the variables should be uniquely determined.

The only way $\lambda_1$ is uniquely determined is if the inequality binds: $\lambda_1 = p_1 - p_2$.

If $\lambda_1$ could be anything in $[0, p_1 - p_2]$, we wouldn't have a unique solution to the KKT system, which contradicts the fact that our original optimization problem has a unique solution.

\begin{tcolorbox}[colback=green!5!white,colframe=green!75!black,title=Uniqueness Argument]
\textbf{The original problem} (minimize cost subject to constraints) has a unique solution because perfect substitutes always lead to corner solutions.

\textbf{Therefore}, the KKT system must also have a unique solution.

\textbf{Therefore}, all variables (including $\lambda_1$) must be uniquely determined.

\textbf{Therefore}, the inequality must bind: $\lambda_1 = p_1 - p_2$.
\end{tcolorbox}

\end{document}
